% STATUS: written 2026-11-05 (Alaminos--Bresar--Extremera--Villena 2009 read in
% full; excerpt in excerpts/. The 2009 paper's method extends the
% Chebotar--Ke--Lee approach recalled in Section 03.)
\section{Derivations among preserver problems}\label{sec:local-characterizations:preservers}

The zero product condition of Section~\ref{sec:local-characterizations:zero-products}
came from a preserver question. A linear map $T$ between algebras
\emph{preserves zero products} when $T(a)T(b)=0$ whenever $ab=0$
\cite[p.~131]{AlaminosBresar2009}. Banach started this line: a linear map of
$L_p([0,1])$ that keeps distances sends two functions whose supports do not
meet to two functions whose supports do not meet
\cite[p.~131]{AlaminosBresar2009}. The standard goal
is now to write such a $T$ as a weighted homomorphism
\cite[p.~132]{AlaminosBresar2009}.

Alaminos, Bre\v{s}ar, Extremera and Villena run the same question for
Banach algebras \cite[p.~131]{AlaminosBresar2009}. Their tool is the
bilinear map $\varphi(a,b)=T(a)T(b)$, which vanishes on zero products
exactly when $T$ preserves them \cite[p.~132]{AlaminosBresar2009}.

\begin{theorem}[{\cite[Theorem~2.11, p.~143]{AlaminosBresar2009}}]\label{thm:local:ab211}
Let $A$ be a faithful Banach algebra, let $X$ be a Banach space, and let
$\varphi\colon A\times A\to X$ be a continuous bilinear map preserving zero
products. Then $\varphi(a\mu,b)=\varphi(a,\mu b)$ for all $a,b\in A$ and
$\mu\in\overline{D(A)}^{\,so}$.
\end{theorem}

Here $D(A)$ is the subalgebra of the multiplier algebra $M(A)$ generated by
doubly power-bounded elements, and the closure is taken in the strong
operator topology \cite[p.~143]{AlaminosBresar2009}. The rule of
Theorem~\ref{thm:local:ab211} applied to $\varphi(a,b)=T(a)T(b)$ turns a
surjective zero product preserving map into a weighted homomorphism on a
large class of Banach algebras, group algebras among them
\cite[Theorem~3.6, p.~152]{AlaminosBresar2009}.

The same tool answers the derivational side of
Section~\ref{sec:local-characterizations:zero-products} for Banach algebras
\cite[p.~155]{AlaminosBresar2009}. A continuous linear $D\colon A\to X$
with $D(a)\cdot b+a\cdot D(b)=0$ whenever $ab=0$ need not be a derivation:
by Theorem~\ref{thm:local:ab46} it lands in a wider class.

\begin{theorem}[{\cite[Theorem~4.6, p.~156]{AlaminosBresar2009}}]\label{thm:local:ab46}
Let $A$ be a Banach algebra with the property of
Theorem~\ref{thm:local:ab211} and a bounded approximate identity, let $X$
be a Banach $A$-bimodule with $\operatorname{Ann}_X(A)=\{0\}$, and let
$D\colon A\to X$ be a continuous linear operator satisfying
$D(a)\cdot b+a\cdot D(b)=0$ whenever $ab=0$. Then $D$ is a generalized
derivation: there exists $\xi\in Z_A(X^{**})$ such that
$D(ab)=D(a)\cdot b+a\cdot D(b)-a\cdot\xi\cdot b$ for all $a,b\in A$.
Accordingly, there exists a continuous derivation $d\colon A\to X^{**}$
such that $D(a)=d(a)+a\cdot\xi$ for all $a\in A$. Moreover $\xi$ can be
chosen in $Z_A(X)$ when $A$ has an identity or $X$ is a dual bimodule.
\end{theorem}

Theorem~\ref{thm:local:ab46} is the Banach cousin of the prime-ring
Theorem~\ref{thm:local:cklzero}. The correction element $\lambda$ of
Theorem~\ref{thm:local:cklzero} becomes a central $\xi$ in the second dual
\cite[p.~156]{AlaminosBresar2009}. When $A$ is
unital, $\xi$ sits in the bimodule itself
\cite[p.~156]{AlaminosBresar2009}.
